\chapter{Exhaustive Astrobiological Conditional Probability Tables (CPTs)}\label{app:complete_cpts}

This appendix provides the full, unabridged mathematical specification and astrophysical justification of the Conditional Probability Tables (CPTs) parameterizing the 9-node Bayesian belief network for the exoplanetary system K2-18b.

\section{Graph Topology and Variable State Spaces}\label{sec:app_graph_topology}

The astrobiological network is formalized as a Directed Acyclic Graph (DAG) $\mathcal{G} = (\mathbf{X}, \mathbf{E})$ over nine discrete binary random variables $\mathbf{X} = \{X_0, X_1, \dots, X_8\}$, where each variable represents a physical, atmospheric, or biological state:
\begin{equation}
X_i \in \{0, 1\}, \quad \forall i \in \{0, 1, \dots, 8\}.
\end{equation}
The topological layout is organized across three hierarchical causal tiers:
\begin{enumerate}[leftmargin=*]
    \item \textbf{Astrophysical Foundation Tiers (Root Nodes):}
    \begin{itemize}
        \item $X_0 \in \{0, 1\}$: Stellar host classification ($X_0 = 1$ denotes an active M-dwarf host).
        \item $X_1 \in \{0, 1\}$: Orbital zone ($X_1 = 1$ denotes a temperate orbit within the conservative liquid water habitable zone).
    \end{itemize}
    \item \textbf{Planetary and Geochemical Interior Tiers:}
    \begin{itemize}
        \item $X_2 \in \{0, 1\}$: Planetary structural regime ($X_2 = 1$ denotes a Hycean sub-Neptune regime characterized by a substantial $\mathrm{H}_2$-rich envelope overlying a global ocean mantle).
        \item $X_3 \in \{0, 1\}$: Marine surface condition ($X_3 = 1$ denotes an active marine biosphere sustaining geochemical recycling).
    \end{itemize}
    \item \textbf{Spectroscopic Transmission and Biosignature Tiers:}
    \begin{itemize}
        \item $X_4 \in \{0, 1\}$: Atmospheric methane absorption at $1.4\,\mu\mathrm{m}$, $2.3\,\mu\mathrm{m}$, and $3.3\,\mu\mathrm{m}$ ($X_4 = 1$ denotes $>5\sigma$ spectral absorption).
        \item $X_5 \in \{0, 1\}$: Atmospheric carbon dioxide absorption at $2.7\,\mu\mathrm{m}$ and $4.3\,\mu\mathrm{m}$ ($X_5 = 1$ denotes $\sim 3\sigma$ spectral absorption).
        \item $X_6 \in \{0, 1\}$: Tropospheric water vapor absorption at $1.4\,\mu\mathrm{m}$ ($X_6 = 1$ denotes detected absorption).
        \item $X_7 \in \{0, 1\}$: Tentative volatile organosulfur biosignature candidate (dimethyl sulfide, $\mathrm{DMS}$, $\mathrm{C}_2\mathrm{H}_6\mathrm{S}$).
        \item $X_8 \in \{0, 1\}$: Industrial synthetic halocarbon technosignature (Freon-12, $\mathrm{CF}_2\mathrm{Cl}_2$).
    \end{itemize}
\end{enumerate}

The joint probability distribution factorizes according to the chain rule of Bayesian networks:
\begin{equation}
P(X_0, X_1, \dots, X_8) = \prod_{i=0}^8 P(X_i \mid \mathrm{Pa}(X_i)),
\label{eq:app_joint_factorization}
\end{equation}
where $\mathrm{Pa}(X_i)$ denotes the set of immediate parents of node $X_i$ in $\mathcal{G}$.

\section{Exhaustive Combinatorial CPT Specifications}\label{sec:app_cpt_specifications}

Tables \ref{tab:cpt_x0} through \ref{tab:cpt_x8} detail every combinatorial state assignment for each conditional distribution, together with the physical mechanism governing the transition probabilities.

\subsection{Node $X_0$: Stellar Host Classification}
\begin{table}[htbp]
\centering
\caption{\textbf{Conditional Probability Table for Node $X_0$ (Stellar M-Dwarf).}}
\label{tab:cpt_x0}
\small
\begin{tabular}{lcc}
\toprule
\textbf{Parents $\mathrm{Pa}(X_0)$} & $P(X_0 = 0)$ & $P(X_0 = 1)$ \\
\midrule
$\emptyset$ (Root Prior) & $0.2500$ & $0.7500$ \\
\bottomrule
\end{tabular}
\end{table}
\textbf{Astrophysical Rationale:} K2-18 (EPIC 201912552) is an established M2.5V dwarf star with mass $M_* = 0.495 \pm 0.004\,M_\odot$, radius $R_* = 0.469 \pm 0.010\,R_\odot$, and effective temperature $T_{\mathrm{eff}} = 3457 \pm 39\,\mathrm{K}$ \cite{benneke2019}. We assign a high prior baseline of $0.75$ reflecting the observational certainty of its low-mass stellar classification while accounting for uncertainties in chromospheric activity and flaring frequency.

\subsection{Node $X_1$: Orbital Habitable Zone}
\begin{table}[htbp]
\centering
\caption{\textbf{Conditional Probability Table for Node $X_1$ (Habitable Zone Orbit).}}
\label{tab:cpt_x1}
\small
\begin{tabular}{lcc}
\toprule
\textbf{Parents $\mathrm{Pa}(X_1)$} & $P(X_1 = 0)$ & $P(X_1 = 1)$ \\
\midrule
$\emptyset$ (Root Prior) & $0.8000$ & $0.2000$ \\
\bottomrule
\end{tabular}
\end{table}
\textbf{Astrophysical Rationale:} K2-18b orbits at a semi-major axis $a = 0.1429 \pm 0.0006\,\mathrm{AU}$ with an orbital period of $32.9396\,\mathrm{days}$. The incident stellar flux is $S = 1.22 \pm 0.08\,S_\oplus$ ($1368\,\mathrm{W/m^2} \times 0.94$), placing the planet near the inner edge of the conservative liquid water habitable zone for an M-dwarf host \cite{benneke2019, madhusudhan2023}. A prior of $P(X_1 = 1) = 0.20$ reflects the geometric and albedo margin required to avoid a runaway greenhouse state.

\subsection{Node $X_2$: Hycean Planetary Structural Regime}
\begin{table}[htbp]
\centering
\caption{\textbf{Conditional Probability Table for Node $X_2$ (Hycean Planet Regime).}}
\label{tab:cpt_x2}
\small
\begin{tabular}{cccc}
\toprule
\textbf{Host Star ($X_0$)} & \textbf{Habitable Orbit ($X_1$)} & $P(X_2 = 0 \mid X_0, X_1)$ & $P(X_2 = 1 \mid X_0, X_1)$ \\
\midrule
$0$ & $0$ & $0.9500$ & $0.0500$ \\
$0$ & $1$ & $0.9500$ & $0.0500$ \\
$1$ & $0$ & $0.9500$ & $0.0500$ \\
$1$ & $1$ & $0.1500$ & $0.8500$ \\
\bottomrule
\end{tabular}
\end{table}
\textbf{Geophysical Rationale:} Hycean worlds require a low-density volatile envelope ($R_p = 2.61 \pm 0.09\,R_\oplus$, $M_p = 8.63 \pm 1.35\,M_\oplus$, yielding bulk density $\rho \approx 2.67\,\mathrm{g/cm^3}$) overlying a liquid ocean \cite{madhusudhan2023}. When both an active M-dwarf providing suitable UV photolysis ($X_0=1$) and a temperate insolation orbit ($X_1=1$) coincide, the probability of establishing a stable Hycean equilibrium reaches $0.85$. Any non-habitable or non-M-dwarf configuration collapses the likelihood of a stable Hycean envelope to a residual baseline of $0.05$.

\subsection{Node $X_3$: Active Marine Biosphere}
\begin{table}[htbp]
\centering
\caption{\textbf{Conditional Probability Table for Node $X_3$ (Marine Biosphere).}}
\label{tab:cpt_x3}
\small
\begin{tabular}{ccc}
\toprule
\textbf{Hycean Regime ($X_2$)} & $P(X_3 = 0 \mid X_2)$ & $P(X_3 = 1 \mid X_2)$ \\
\midrule
$0$ & $0.9900$ & $0.0100$ \\
$1$ & $0.6000$ & $0.4000$ \\
\bottomrule
\end{tabular}
\end{table}
\textbf{Astrobiological Rationale:} Under a confirmed Hycean ocean with a habitable surface boundary layer ($X_2=1$), hydrothermal circulation and chemical disequilibrium provide favorable thermodynamic gradients for chemosynthetic marine life, parameterized at $P(X_3=1 \mid X_2=1) = 0.40$. Without a Hycean ocean ($X_2=0$), the likelihood of an active biosphere drops to an abiotic background of $0.01$.

\subsection{Node $X_4$: Transmission Feature of Methane ($\mathrm{CH}_4$)}
\begin{table}[htbp]
\centering
\caption{\textbf{Conditional Probability Table for Node $X_4$ ($\mathrm{CH}_4$ Absorption).}}
\label{tab:cpt_x4}
\small
\begin{tabular}{ccc}
\toprule
\textbf{Hycean Regime ($X_2$)} & $P(X_4 = 0 \mid X_2)$ & $P(X_4 = 1 \mid X_2)$ \\
\midrule
$0$ & $0.8500$ & $0.1500$ \\
$1$ & $0.1000$ & $0.9000$ \\
\bottomrule
\end{tabular}
\end{table}
\textbf{Spectroscopic Rationale:} Observations with JWST NIRISS and NIRSpec revealed strong methane absorption features ($>5\sigma$ detection) in K2-18b's transmission spectrum \cite{madhusudhan2023}. In a reducing $\mathrm{H}_2$-rich Hycean atmosphere, methane is the dominant equilibrium carbon-bearing molecule, yielding $P(X_4=1 \mid X_2=1) = 0.90$. In a non-Hycean regime (e.g., bare rocky core or mini-Neptune with high metallicity), methane is readily depleted by photolysis, parameterized at $0.15$.

\subsection{Node $X_5$: Transmission Feature of Carbon Dioxide ($\mathrm{CO}_2$)}
\begin{table}[htbp]
\centering
\caption{\textbf{Conditional Probability Table for Node $X_5$ ($\mathrm{CO}_2$ Absorption).}}
\label{tab:cpt_x5}
\small
\begin{tabular}{ccc}
\toprule
\textbf{Hycean Regime ($X_2$)} & $P(X_5 = 0 \mid X_2)$ & $P(X_5 = 1 \mid X_2)$ \\
\midrule
$0$ & $0.8000$ & $0.2000$ \\
$1$ & $0.2000$ & $0.8000$ \\
\bottomrule
\end{tabular}
\end{table}
\textbf{Spectroscopic Rationale:} JWST transmission observations detected $\mathrm{CO}_2$ at $\sim 3\sigma$ confidence \cite{madhusudhan2023}. Carbon dioxide production in a temperate Hycean atmosphere is sustained by photochemical oxidation of carbon monoxide and ocean-atmosphere volatile exchange ($P = 0.80$). Without a Hycean ocean to act as a chemical buffer, $\mathrm{CO}_2$ occurrence is parameterized at $0.20$.

\subsection{Node $X_6$: Tropospheric Water Vapor ($\mathrm{H}_2\mathrm{O}$)}
\begin{table}[htbp]
\centering
\caption{\textbf{Conditional Probability Table for Node $X_6$ ($\mathrm{H}_2\mathrm{O}$ Vapor).}}
\label{tab:cpt_x6}
\small
\begin{tabular}{ccc}
\toprule
\textbf{Habitable Orbit ($X_1$)} & $P(X_6 = 0 \mid X_1)$ & $P(X_6 = 1 \mid X_1)$ \\
\midrule
$0$ & $0.9000$ & $0.1000$ \\
$1$ & $0.3000$ & $0.7000$ \\
\bottomrule
\end{tabular}
\end{table}
\textbf{Spectroscopic Rationale:} Water vapor in the lower troposphere depends directly on orbital insolation ($X_1$). Within the habitable zone ($X_1=1$), liquid surface evaporation maintains atmospheric water vapor ($P = 0.70$), as observed historically by HST \cite{benneke2019}. Outside the habitable zone ($X_1=0$), cold-trapping or runaway photolysis suppresses gaseous $\mathrm{H}_2\mathrm{O}$ detection down to $0.10$.

\subsection{Node $X_7$: Dimethyl Sulfide ($\mathrm{DMS}$) Biosignature}
\begin{table}[htbp]
\centering
\caption{\textbf{Conditional Probability Table for Node $X_7$ ($\mathrm{DMS}$ Biosignature).}}
\label{tab:cpt_x7}
\small
\begin{tabular}{ccc}
\toprule
\textbf{Marine Biosphere ($X_3$)} & $P(X_7 = 0 \mid X_3)$ & $P(X_7 = 1 \mid X_3)$ \\
\midrule
$0$ & $0.9999$ & $0.0001$ \\
$1$ & $0.9500$ & $0.0500$ \\
\bottomrule
\end{tabular}
\end{table}
\textbf{Biogeochemical Rationale:} Dimethyl sulfide ($\mathrm{CH}_3\mathrm{SCH}_3$) is produced on Earth by marine organisms through the enzymatic cleavage of dimethylsulfoniopropionate (DMSP). In an active marine biosphere ($X_3=1$), DMS can accumulate to detectable levels ($P = 0.05$). In the absence of biology ($X_3=0$), abiotic photochemical synthesis pathways are exceedingly inefficient, yielding an infinitesimal false-positive rate of $10^{-4}$ ($0.01\%$). While tentatively reported in K2-18b by Madhusudhan et al.~\cite{madhusudhan2023}, this feature remains subject to ongoing debate in the astrobiological community and requires confirmation through future high-precision spectroscopic passes.

\subsection{Node $X_8$: Industrial Chlorofluorocarbon ($\mathrm{CFC}$) Technosignature}
\begin{table}[htbp]
\centering
\caption{\textbf{Conditional Probability Table for Node $X_8$ ($\mathrm{CFC}$ Technosignature).}}
\label{tab:cpt_x8}
\small
\begin{tabular}{ccc}
\toprule
\textbf{Marine Biosphere ($X_3$)} & $P(X_8 = 0 \mid X_3)$ & $P(X_8 = 1 \mid X_3)$ \\
\midrule
$0$ & $1.00000000$ & $0.00000000$ \\
$1$ & $0.99998689$ & $1.310616 \times 10^{-5}$ \\
\bottomrule
\end{tabular}
\end{table}
\textbf{Technosignature Rationale:} Chlorofluorocarbons (specifically $\mathrm{CF}_2\mathrm{Cl}_2$, CFC-12) have no known natural abiotic production pathways in planetary atmospheres \cite{lin2014, haqqmisra2022}. Consequently, $P(X_8 = 1 \mid X_3 = 0) = 0$ identically. Conditioned on an active biosphere capable of technological development ($X_3=1$), the transition probability is calibrated to enforce the exact canonical prior:
\begin{equation}
P(X_8 = 1 \mid X_3 = 1) = \frac{p_{\mathrm{rare}}}{P(X_3 = 1)} = \frac{1.0 \times 10^{-6}}{0.07630} = 1.310616 \times 10^{-5}.
\end{equation}

\section{Exact Marginal and Conditional Posteriors}\label{sec:app_exact_posteriors}

Evaluating the marginal distributions analytically through the DAG topology yields:
\begin{align}
P(X_0 = 1) &= 0.750000, \\
P(X_1 = 1) &= 0.200000, \\
P(X_2 = 1) &= 0.85(0.75 \times 0.20) + 0.05(1 - 0.75 \times 0.20) = 0.170000, \\
P(X_3 = 1) &= 0.40(0.1700) + 0.01(0.8300) = \mathbf{0.076300}, \\
P(X_4 = 1) &= 0.90(0.1700) + 0.15(0.8300) = 0.277500, \\
P(X_5 = 1) &= 0.80(0.1700) + 0.20(0.8300) = 0.302000, \\
P(X_6 = 1) &= 0.70(0.2000) + 0.10(0.8000) = 0.220000, \\
P(X_7 = 1) &= 0.05(0.0763) + 0.0001(0.9237) = 0.003907, \\
P(X_8 = 1) &= (1.310616 \times 10^{-5})(0.0763) + 0 = \mathbf{1.000000 \times 10^{-6}}.
\end{align}

Under transmission evidence matching the JWST observation $\mathbf{e} = \{X_4 = 1, X_5 = 1, X_6 = 1\}$:
\begin{align}
P(\mathbf{e}) &= \sum_{\mathbf{x} \setminus \{X_4, X_5, X_6\}} P(\mathbf{x}, \mathbf{e}) = \mathbf{0.095804} \quad (9.58\%), \\
P(X_7 = 1 \mid \mathbf{e}) &= \frac{P(X_7 = 1, \mathbf{e})}{P(\mathbf{e})} = \mathbf{0.019630} \quad (1.96\%), \\
P(X_8 = 1 \mid \mathbf{e}) &= \frac{P(X_8 = 1, \mathbf{e})}{P(\mathbf{e})} = \mathbf{5.130310 \times 10^{-6}}.
\end{align}
These exact analytical values represent the ground-truth benchmark evaluated in Chapter \ref{chap:classical_limits} and Chapter \ref{chap:system_architecture}.

\subsection{Methodological Invariance of Algorithmic Complexity}\label{subsec:app_algorithmic_invariance}
It is essential to emphasize that the foundational conclusions of this thesis—specifically, the exponential spatial compression $\mathcal{O}(N)$ achieved via Amplitude Encoding and the quadratic Heisenberg speedup $\mathcal{O}(1/M)$ demonstrated by QAE—are intrinsic mathematical properties of the quantum circuit topology and unitary reflection dynamics. While the numerical probabilities defined in Tables \ref{tab:cpt_x0}--\ref{tab:cpt_x8} reflect physically grounded astrophysical and photochemical priors for K2-18b, the algorithmic speedups and variance stabilization over classical MCMC are strictly invariant to the specific scalar entries of the CPTs, holding universally for any sparse causal Bayesian network.
